%% D-P2-010: elementary semantic results for the active source CPL language.
\subsection{Uniform substitution and replacement}
\label{sec:p02-propositional-substitution}

Use the alphabet and formula strings of
Section~\ref{sec:prop-syntax-semantics-01} and the valuations of
Section~\ref{sec:prop-syntax-semantics-02}.  Thus
\(\mathsf{Prop}=\{p_n:n\in\mathbb N\}\), formulas belong to
\(\mathsf{Fm}_{\mathsf{CPL}}\), and the primitive constructors are exactly
\(\bottom,\neg,\land,\lor,\to\).  The two truth values are
\(\mathbf F,\mathbf T\), with the recursive extension of \(v\) written
\(\widehat v\).  In this presentation
\(\alpha\equiv\beta\) retains the preceding semantic meaning: every valuation
gives \(\alpha\) and \(\beta\) the same truth value.  It is a metalinguistic
equivalence relation, not an additional formula constructor.  Defined
connectives such as \(\topv\) and \(\leftrightarrow\) are expanded into the
primitive grammar before the operations below are applied.

\begin{hbdefinition}{def:p02-uniform-substitution}{Uniform substitution}
A \emph{uniform substitution} is a function
\(\upsilon:\mathsf{Prop}\to\mathsf{Fm}_{\mathsf{CPL}}\).  Its simultaneous
extension to formula strings is defined by
\[
\begin{aligned}
 p_n[\upsilon]&=\upsilon(p_n),&
 \bottom[\upsilon]&=\bottom,\\
 (\neg\alpha)[\upsilon]&=\neg(\alpha[\upsilon]),\\
 (\alpha\land\beta)[\upsilon]&=(\alpha[\upsilon]\land\beta[\upsilon]),\\
 (\alpha\lor\beta)[\upsilon]&=(\alpha[\upsilon]\lor\beta[\upsilon]),\\
 (\alpha\to\beta)[\upsilon]&=(\alpha[\upsilon]\to\beta[\upsilon]).
\end{aligned}
\]
At an original variable occurrence, insert its assigned formula once.
Do not apply \(\upsilon\) again inside an inserted formula.
Theorem~\ref{thm:prop-unique-readability} makes this a well-defined structural
recursion, and every displayed output is a formula of the same language.

For a valuation \(v\), define another valuation
\(v_\upsilon:\mathsf{Prop}\to\{\mathbf F,\mathbf T\}\) by
\[
 v_\upsilon(p_n)=\widehat v\bigl(\upsilon(p_n)\bigr).
\]
This is defined independently at each variable; no surjectivity of the map
\(v\mapsto v_\upsilon\) is assumed.
\index{uniform substitution!propositional}
\end{hbdefinition}

\begin{hbtheorem}{thm:p02-uniform-substitution}{proved=annex:p2.prop.substitution}{Semantic uniform-substitution theorem}
For every uniform substitution \(\upsilon\), valuation \(v\), and formulas
\(\varphi,\psi\),
\[
 \widehat v\bigl(\varphi[\upsilon]\bigr)
   =\widehat{v_\upsilon}(\varphi).
\]
Consequently,
\[
\begin{aligned}
 \entails\varphi&\quad\Longrightarrow\quad
       \entails\varphi[\upsilon],\\
 \varphi\equiv\psi&\quad\Longrightarrow\quad
       \varphi[\upsilon]\equiv\psi[\upsilon].
\end{aligned}
\]
\end{hbtheorem}

\begin{hbdefinition}{def:p02-selected-replacement}{Selected subformula occurrences and replacement}
An occurrence is a position in the original formula's parse tree.
Represent positions by finite words over \(\{1,2\}\), with the empty word
\(\epsilon\) denoting the root.  If \(i\) is a digit and \(a\) is a position,
\(ia\) denotes their concatenation.  The finite position sets are
\[
\begin{aligned}
 \operatorname{Pos}(p_n)&=\operatorname{Pos}(\bottom)=\{\epsilon\},\\
 \operatorname{Pos}(\neg\alpha)
   &=\{\epsilon\}\cup\{1a:a\in\operatorname{Pos}(\alpha)\},\\
 \operatorname{Pos}(\alpha\circ\beta)
  &=\{\epsilon\}\cup\{1a:a\in\operatorname{Pos}(\alpha)\}\\
  &\quad\cup\{2a:a\in\operatorname{Pos}(\beta)\}.
\end{aligned}
\]
Here \(\circ\in\{\land,\lor,\to\}\).
Write \(\varphi_a\) for the subformula at position \(a\):
\(\varphi_\epsilon=\varphi\),
\((\neg\alpha)_{1a}=\alpha_a\), and
\((\alpha\circ\beta)_{1a}=\alpha_a\),
\((\alpha\circ\beta)_{2a}=\beta_a\) whenever the position exists.

Fix formulas \(\varphi,\psi,\chi\) and a set
\[
 O\subseteq\{a\in\operatorname{Pos}(\varphi):\varphi_a=\psi\}.
\]
These are the chosen occurrences of the formula string \(\psi\), not merely
occurrences of its variables.  No chosen position is a proper prefix of
another: a proper subformula has strictly fewer occurrences of
\(\neg,\land,\lor,\to\), so a finite formula cannot be a proper subformula
of itself.
The set \(O\) may be empty or may contain every occurrence of \(\psi\).

The replacement \(\varphi[O\leftarrow\chi]\) is obtained from the original
tree as follows.  If \(\epsilon\in O\), the result is \(\chi\).  Otherwise put
\(O_i=\{a:ia\in O\}\) and use the recursion
\[
\begin{aligned}
 p_n[\varnothing\leftarrow\chi]&=p_n,&
 \bottom[\varnothing\leftarrow\chi]&=\bottom,\\
 (\neg\alpha)[O\leftarrow\chi]
   &=\neg\bigl(\alpha[O_1\leftarrow\chi]\bigr),\\
 (\alpha\circ\beta)[O\leftarrow\chi]
   &=\bigl(\alpha[O_1\leftarrow\chi]\circ
           \beta[O_2\leftarrow\chi]\bigr),
 \qquad \circ\in\{\land,\lor,\to\}.
\end{aligned}
\]
Thus all selected original subtrees are replaced simultaneously; no inserted
copy of \(\chi\) is searched for further replacements.
\index{replacement!equivalent subformulas}
\index{subformula!occurrences}
\end{hbdefinition}

\begin{hbtheorem}{thm:p02-replacement}{proved=annex:p2.prop.replacement}{Equivalence congruence and selected-occurrence replacement}
Logical equivalence is reflexive, symmetric, and transitive.  It is a
congruence for every primitive connective: \(\bottom\equiv\bottom\), and
\[
 \alpha\equiv\alpha'
   \quad\Longrightarrow\quad \neg\alpha\equiv\neg\alpha'.
\]
If \(\alpha\equiv\alpha'\) and \(\beta\equiv\beta'\), then
\[
\begin{aligned}
 (\alpha\land\beta)&\equiv(\alpha'\land\beta'),\\
 (\alpha\lor\beta)&\equiv(\alpha'\lor\beta'),\\
 (\alpha\to\beta)&\equiv(\alpha'\to\beta').
\end{aligned}
\]
If \(\psi\equiv\chi\) and \(O\) is any chosen set of occurrences of
\(\psi\) in \(\varphi\), then
\[
 \varphi\equiv\varphi[O\leftarrow\chi].
\]
\end{hbtheorem}

\begin{hbexample}{def:p02-substitution-replacement-example}{Simultaneous substitution versus one selected replacement}
Fix distinct propositional variables \(P,Q\).  Let
\(\upsilon(P)=Q\), \(\upsilon(Q)=\neg Q\), and let every other variable be
fixed.  Then
\[
\begin{aligned}
 (P\land P)[\upsilon]&=(Q\land Q),\\
 (P\land P)[\{1\}\leftarrow Q]&=(Q\land P).
\end{aligned}
\]
In the first expression the inserted \(Q\)'s are not substituted again.
In the second expression only the left occurrence of \(P\) is selected.
The two outputs differ in truth value when \(v(P)=\mathbf F\) and
\(v(Q)=\mathbf T\).  This comparison concerns the operations themselves;
\(P\) and \(Q\) are not logically equivalent, so the replacement theorem does
not assert equivalence with the original formula in this example.

For an equivalence-preserving application, put
\(\psi=(P\to Q)\) and \(\chi=(\neg P\lor Q)\).
The truth table in Section~\ref{sec:prop-syntax-semantics-02} gives
\(\psi\equiv\chi\).  Replacing only the left occurrence of \(\psi\) in
\((\psi\land\psi)\) therefore yields the equivalent formula
\((\chi\land\psi)\); the right occurrence remains unchanged.
\end{hbexample}

\begin{hbexample}{def:p02-substitution-replacement-limits}{Two invalid weakenings of the hypotheses}
Validity of one substitution instance need not imply validity of the original.
Take \(\varphi=P\) and \(\upsilon(P)=\neg\bottom\), fixing all other
variables.  The instance is valid, whereas \(P\) is false under a valuation
with \(v(P)=\mathbf F\).

Merely equisatisfiable formulas cannot replace equivalent formulas in the
replacement theorem.  Here \emph{equisatisfiable} means that one formula has a
satisfying valuation if and only if the other does.  Take distinct variables
\(P,Q\), which are both satisfiable.  Nevertheless, replacing only the left \(P\) in
\((P\land\neg P)\) by \(Q\) gives \((Q\land\neg P)\).
The original is false under every valuation; the result is true when
\(v(P)=\mathbf F\) and \(v(Q)=\mathbf T\).  This replacement preserves neither
logical equivalence nor even satisfiability of the surrounding formula.
\end{hbexample}
